\begin{theorem}[Conductors under trace and norm]
\label{mod:ramification-pullbackconductors}
Let $K/F$ be a finite Galois valuative extension of nonarchimedean local
fields.  Fix an integral uniformizer $\pi_K$ such that
\[
 \operatorname{Algebra.adjoin}_{\mathcal O_F}\{\pi_K\}=\mathcal O_K,
\]
and put
\[
 e=e(K/F),\qquad D=d_{K/F}=v_K(\mathfrak D_{K/F}).
\]
If $\psi$ has exact additive conductor $n\in\mathbf Z$, in the convention
that $\mathfrak p_F^{-n}$ is its largest triviality lattice, then
\[
 \boxed{\quad
 n_K(\psi\circ\operatorname{Tr}_{K/F})=e n+D.
 \quad}
\]
Thus the different shift has a plus sign in the conductor and a minus sign
in the largest trivial lattice.  In particular, the statement is valid for
arbitrary positive, zero, or negative $n$.

Now assume that $K/F$ is cyclic of prime degree, that $t\in\mathbf N$ is a
supplied lower-break witness, and that $f(K/F)=1$.  Retain the chosen
generating uniformizer and write
\[
 \ell=[K:F],\qquad D=(\ell-1)(t+1),\qquad
 \Psi(r)=\Psi_{t,\ell}(r).
\]
Then $e=\ell$, so the additive formula specializes to
\[
 \boxed{\quad
 n_K(\psi\circ\operatorname{Tr}_{K/F})
 =\ell n_F(\psi)+(\ell-1)(t+1).
 \quad}
\]

Let $\chi$ be a continuous quasi-character of $F^\times$, and suppose its
exact conductor is $m$.  Write $m'$ for the exact conductor of
$\chi\circ N_{K/F}$.  Norm-pullback triviality on a source layer is first
identified with triviality of $\chi$ on its honest full norm image.  Hence
the exact image equalities from
module~\ref{mod:ramification-normfiltration} give
\[
 \operatorname{MulTriv}_K(\chi\circ N,U_K^{\Psi(r)})
 \Longleftrightarrow
 \operatorname{MulTriv}_F(\chi,U_F^r),
 \qquad t<r,
\]
and, crucially, the separate adjacent-denominator equivalence
\[
 \operatorname{MulTriv}_K(\chi\circ N,U_K^{\Psi(r)+1})
 \Longleftrightarrow
 \operatorname{MulTriv}_F(\chi,U_F^{r+1}).
\]
At the critical successor one similarly has
\[
 \operatorname{MulTriv}_K(\chi\circ N,U_K^{t+1})
 \Longleftrightarrow
 \operatorname{MulTriv}_F(\chi,U_F^{t+1}).
\]

These equivalences prove the exact conductor trichotomy, with conductor zero
kept separate from the positive-boundary argument:
\begin{enumerate}[label=(\roman*)]
\item If $m=0$, then $m'=0$.
\item If $1\le m\le t$, then $m'=m$.  Together with the preceding endpoint,
this is the formula $m'=m$ for every $m\le t$.
\item If $m>t+1$, then
\[
 \boxed{\quad m'-1=\Psi(m-1),\qquad
 m'=\Psi(m-1)+1.\quad}
\]
The source depths used here are exactly $\Psi(m-1)$ and
$\Psi(m-1)+1$, never $\Psi(m)$.  Using the prime-cyclic formula for $\Psi$,
this is equivalently
\[
 \boxed{\quad
 m'=\ell m-(\ell-1)(t+1)=\ell m-D.
 \quad}
\]
\end{enumerate}
For the positive below-break case, exactness is not deduced from false full
surjectivity of $U_K^{m-1}\to U_F^{m-1}$.  Instead Lean uses the proved join
identity
\[
 N(U_K^{m-1})\vee U_F^t=U_F^{m-1}
\]
together with triviality on $U_F^m\supseteq U_F^t$.

At the critical conductor $m=t+1$, exact successor surjectivity gives only
the first conclusion
\[
 m'\le t+1.
\]
Lean then proves the cancellation endpoint exactly.  Let
\[
 I_t=N(U_K^t)\le F^\times;
\]
equivalently, $I_t$ is the inverse image in $U_F^t$ of the range of the
critical graded norm.  Then
\[
 \boxed{\quad
 m'<t+1
 \Longleftrightarrow
 (\chi\circ N)|_{U_K^t}=1
 \Longleftrightarrow
 \chi|_{I_t}=1.
 \quad}
\]
Thus failure of this condition proves the exact equality $m'=t+1$, not just
an upper bound.

Write
\[
 S(K/F)=\{\mu:F^\times\to\C^\times:\mu\circ N_{K/F}=1\}.
\]
Without using any theorem from the later norm-character node, Lean also
proves the manuscript's equivalent criterion
\[
 \boxed{\quad
 m'<m=t+1
 \Longleftrightarrow
 \exists\mu\in S(K/F),\ \exists q<t+1,
 \quad m_F(\mu\chi)=q.
 \quad}
\]
The forward implication is constructive at the level of quotient classes.
Triviality on $U_K^t$ factors $\chi\circ N$ through
$K^\times/U_K^t$.  Lean transports the resulting homomorphism through the
proved equivalence
\[
 K^\times/U_K^t\xrightarrow{\ \sim\ }F^\times/U_F^t,
\]
pulls it back to a character $\rho$ of $F^\times$, proves continuity from
its open unit-layer kernel, and sets $\mu=\rho\chi^{-1}$.  The least unit
layer on which $\rho$ is trivial supplies an exact conductor $q\le t$.
Conversely, $(\mu\chi)\circ N=\chi\circ N$, and the below-break formula
proves the stronger exact conclusion
\[
 m_K(\chi\circ N)=q.
\]
No conductor formula is assumed as a structure field or helper hypothesis.

For packaged exact-conductor data, Lean exports the same equalities as
equalities of the stored conductors.  In particular, if $\chi$ has minimal
conductor in its $S(K/F)$-orbit, then
\[
 m\le t+1\Longrightarrow m'=m.
\]
At and above the critical endpoint this gives the uniform explicit identity
\[
 m'=\ell m-D.
\]
Combining it with the additive formula gives the denominator-valuation
identity needed downstream:
\[
 \boxed{\quad
 m_K+n_K=\ell(m_F+n_F).
 \quad}
\]
Lean also exposes the last-nontrivial-layer forms used in the later case
split:
\[
 m'-1=m-1\quad(m-1<t),
\]
\[
 m'-1=t\quad(m-1=t\text{ and the orbit representative is minimal}),
\]
and
\[
 m'-1=\Psi(m-1)\quad(t<m-1).
\]
These statements include the tame critical endpoint $t=0,m=1$: cancellation
gives conductor $0$, while a minimal-orbit representative has conductor $1$.

\LeanFile{LanglandsFirstMainLemma/Ramification/PullbackConductors.lean}
\lean{LanglandsFirstMainLemma.quasiCharTrivialOnUnitFiltration_compNorm_iff_image}
\lean{LanglandsFirstMainLemma.quasiCharTrivialOnUnitFiltration_compNorm_iff_of_image_eq}
\lean{LanglandsFirstMainLemma.quasiCharTrivialOnUnitFiltration_compNorm_of_maps}
\lean{LanglandsFirstMainLemma.additiveConductor_compTrace}
\lean{LanglandsFirstMainLemma.additiveConductor_compTrace_cyclicPrime}
\lean{LanglandsFirstMainLemma.quasiCharTrivialOnUnitFiltration_compNorm_iff_aboveBreak}
\lean{LanglandsFirstMainLemma.quasiCharTrivialOnUnitFiltration_compNorm_iff_aboveBreak_succ}
\lean{LanglandsFirstMainLemma.quasiCharTrivialOnUnitFiltration_compNorm_iff_break_succ}
\lean{LanglandsFirstMainLemma.multiplicativeConductor_compNorm_zero}
\lean{LanglandsFirstMainLemma.multiplicativeConductor_compNorm_belowBreak_succ}
\lean{LanglandsFirstMainLemma.multiplicativeConductor_compNorm_belowBreak}
\lean{LanglandsFirstMainLemma.multiplicativeConductor_compNorm_aboveBreak}
\lean{LanglandsFirstMainLemma.multiplicativeConductor_compNorm_high}
\lean{LanglandsFirstMainLemma.highConductor_cast_eq_degree_mul_sub_different}
\lean{LanglandsFirstMainLemma.multiplicativeConductor_compNorm_critical_le}
\lean{LanglandsFirstMainLemma.multiplicativeConductor_compNorm_critical_drop_iff}
\lean{LanglandsFirstMainLemma.multiplicativeConductor_compNorm_critical_drop_iff_image}
\lean{LanglandsFirstMainLemma.multiplicativeConductor_compNorm_critical_of_noCancellation}
\lean{LanglandsFirstMainLemma.normCharacter_mul_compNorm}
\lean{LanglandsFirstMainLemma.multiplicativeConductor_compNorm_critical_of_twist}
\lean{LanglandsFirstMainLemma.multiplicativeConductor_compNorm_critical_drop_iff_exists_twist}
\lean{LanglandsFirstMainLemma.LocalAddCharData.conductor_compTrace_eq_cyclicPrime}
\lean{LanglandsFirstMainLemma.LocalQuasiCharData.conductor_compNorm_eq_belowBreak}
\lean{LanglandsFirstMainLemma.LocalQuasiCharData.conductor_compNorm_eq_high}
\lean{LanglandsFirstMainLemma.LocalQuasiCharData.conductor_compNorm_eq_high_explicit}
\lean{LanglandsFirstMainLemma.high_compNorm_add_compTrace_conductor_eq}
\lean{LanglandsFirstMainLemma.LocalQuasiCharData.conductor_compNorm_le_critical}
\lean{LanglandsFirstMainLemma.LocalQuasiCharData.conductor_compNorm_critical_drop_iff_exists_twist}
\lean{LanglandsFirstMainLemma.LocalQuasiCharData.conductor_compNorm_eq_of_minimalOrbit}
\lean{LanglandsFirstMainLemma.LocalQuasiCharData.conductor_compNorm_eq_atOrAboveCritical_of_minimalOrbit}
\lean{LanglandsFirstMainLemma.compNorm_add_compTrace_conductor_eq_atOrAboveCritical_of_minimalOrbit}
\lean{LanglandsFirstMainLemma.LocalQuasiCharData.conductor_compNorm_sub_one_eq_belowBreak}
\lean{LanglandsFirstMainLemma.LocalQuasiCharData.conductor_compNorm_sub_one_eq_critical_of_minimalOrbit}
\lean{LanglandsFirstMainLemma.LocalQuasiCharData.conductor_compNorm_sub_one_eq_aboveBreak}
\leanok
\SourceText{Corollary~\ref{cor:additive-conductor} and
Proposition~\ref{prop:pullback-conductor}.}
\uses{mod:ramification-normfiltration,mod:basic-characterconductors}
\FormalizationContract
The additive convention is fixed by the largest trivial lattice
$\mathfrak p^{-n}$, so the different appears with a plus sign in the
conductor.  The different is normalized upstairs.  The lower break $t$ and
inverse Herbrand direction are retained exactly.  Above the break, the two
source layers are $\Psi(r)$ and $\Psi(r)+1$; the proof never substitutes
$\Psi(r+1)$ for the latter.

Conductor zero is proved without a predecessor layer.  Positive conductors
use adjacent exact boundary conditions.  The cases $m\le t$, $m=t+1$, and
$t+1<m$ are disjoint and exhaustive, and the endpoints $t=0,m=1$ and
$m=t+2$ are not conflated.  Below the break the proof uses the exact join
description of the norm image, at the critical break it uses the honest
critical image and exact successor image, and above the break it uses both
exact Herbrand-indexed full images.

The critical norm character is constructed from quotient maps and their
proved equivalence.  Quotient well-definedness, the open-kernel continuity
argument, the exact least conductor of the constructed twist, and equality
of its norm pullback are explicit.  The proof imports no later parameter,
case, norm-character conductor, or local-class-field-theory result.
\end{theorem}
